\section{Local cable tests and the no-crossing principle}
\label{sec:local-kernel}

A canonical cluster is a function of the limiting field, whereas the
stopped law in \cref{prop:exploration-thin} is conditioned on retained
exploration data. To compare them, we turn a local cable path test into
a bounded function of the continuum field on a slightly larger buffer.
Calibration in the canonical parent makes this function equal to one
on a crossing event. Its expectation in the stopped component tends to
zero under that component's conditional GFF law.

\subsection{Auxiliary cables and the tests}

All cable fields in this section are auxiliary Gaussian approximations;
their index \(n\) and dyadic mesh \(\varepsilon_n=2^{-n}\) are separate from
the current mesh \(\delta\) in \cref{sec:local-sets}. They have the
normalization of \(h\), with Cameron--Martin norm
\begin{equation}
 \|u\|_{\nabla}^{2}=\frac1{2\pi}\int |\nabla u|^{2}\,\dd z,
 \qquad u\in H^1_0.
 \label{eq:kernel-energy}
\end{equation}
Fix the corresponding square cable grids \(\Gamma_n\) in one bounded
box \(\mathcal B\). All later subsequences are drawn from this dyadic
sequence. A cable GFF
\(\widetilde\Phi_n\) has discrete GFF values at vertices and independent
Brownian bridges on the metric edges conditional on those values. Edge
resistances and the field scaling are fixed once for this normalization.
For planar pairings we always use the ALS harmonic extension of the
\emph{whole cable function} to the box, with zero data at the box boundary
\cite[Sections~2.1 and~4.1]{ALS2018}; denote it by
\(\widehat\Phi_n\). Thus
\(\widehat\Phi_n(f)=\int f(z)\widehat\Phi_n(z)\,\dd z\).
It is not a vertex-only pairing.
Full-field convergence and full-field functionals below use
\(H^s(\mathcal B)\) for one fixed \(s<-1\), the distribution topology
of the cited cable convergence.

For every test polygon \(B\), fix once a connected closed union
\(B_n^\square\) of complete primal mesh cells approximating
\(\overline B\) from inside and exhausting its compact subsets. Put
\(C_{B,n}=\Gamma_n\cap B_n^\square\). It includes the contained and
perimeter edges, but not the open part of an outward edge. The same
convention defines \(C_{K,n}\) for a smaller test polygon.

The reference polygon \(P\) uses a distinct approximation so that it
is literally a metric domain of \cite[Section~4.1]{ALS2018}. Let
\(P_n^\diamond\) be the interior of a fixed simply connected union of dual
mesh cells centered at grid vertices, exhausting compact subsets of
\(P\) with complement-Hausdorff convergence. Its boundary crosses
primal cables at their midpoints. Subdivide there and write
\(\mathsf P_n=\Gamma_n\cap P_n^\diamond\), including the cut points
as Dirichlet boundary vertices. The inherited half-edge resistances
preserve the cable law. In every compared ambient approximation,
this same subdivided \(\mathsf P_n\) must occur unchanged for fine
mesh, with no missing cable, altered resistance, or extra pin there.
The cut points are unpinned in the ambient field; they become Dirichlet
vertices only in the reference residual after the exterior is conditioned.
The reference is the zero-boundary cable GFF in \(\mathsf P_n\).
An admissible ambient graph is the source metric domain
\(\Gamma_n\cap D_n\) for a Euclidean open \(D_n\subset\mathcal B\),
with its boundary intersection points adjoined as Dirichlet endpoints.
The regular inner approximations used below exhaust compact subsets
and contain the prescribed \(\mathsf P_n\) for all sufficiently large
\(n\). Their complements converge in Hausdorff distance, and their
numbers of complementary components are uniformly bounded. These are
the admissible approximations in the rest of the section.

Here are the tests used later. A collar choice consists of a compact
rational polygon \(Q\) and regular rational polygons
\[
             Q\Subset V\Subset K\Subset B\Subset P.
\]
For \(t\in a\Z\), let \(T_{n,t}^{+}\) be the indicator that, on the
cables in \(C_{K,n}\), a path in \(\{\widetilde\Phi_n\ge t\}\) joins
\(Q\) to \(K\setminus\overline V\). Define \(T_{n,t}^{-}\) with
\(\widetilde\Phi_n\le t\). A canonical inner hit will be in
\(\operatorname{int}Q\), and its outer hit will have positive distance
from \(\overline V\) and from \(\partial K\). These strict hits persist
under Hausdorff convergence. The tests are measurable functions of the
same cables in every domain. More generally, the next two lemmas allow
any \([0,1]\)-valued measurable test supported in \(C_{K,n}\); no
continuity of that test as a function of a limiting distribution is
assumed.

\subsection{The reverse buffer estimate}

For a fixed polygon \(B\), put
\[
 \mathcal G_n(B)=
 \sigma\bigl(\widetilde\Phi_n(x):x\in C_{B,n}\bigr),\qquad
 \mathcal G(B)=\sigma\bigl(\Phi(f):f\in C_c^\infty(B)\bigr).
\]
The first sigma-field contains the full cable restriction, including
the perimeter bridges. By the cable Markov property, the residual after
conditioning is zero on \(C_{B,n}\) and is an independent killed cable
GFF on its complement. Extend this residual by zero on \(C_{B,n}\)
\emph{before applying the original ambient harmonic extension operator}.
Write \(\mathsf K_n^B(f,g)\) for the covariance of the resulting planar
pairings. In particular, the conditioned perimeter cables remain zero
traces in this extension. We do not replace this form by the extension
from a graph in which those cables were deleted.

\begin{lemma}[Reverse buffer estimate]
\label{lem:reverse-buffer}
Let \(K\Subset B\Subset D\), with \(B\) a regular polygon, and let
\(\widetilde\Phi_n\) be zero-boundary cable GFFs in the admissible
approximations of \(D\). Suppose the full extended Green forms converge
to \(G_D\) on smooth tests. Then
\[
                    \mathsf K_n^B(f,g)
             \longrightarrow G_{D\setminus\overline B}(f,g).
\]
If \(T_n\in[0,1]\) is supported in \(C_{K,n}\) and
\((\widehat\Phi_n,T_n)\) converges jointly along a subsequence to
\((\Phi,T)\), then
\begin{equation}
       \E[T\mid\Phi]=\E[T\mid\mathcal G(B)].
 \label{eq:kernel-localization}
\end{equation}
The same conclusion holds for a bounded vector of local tests and after
adding deterministic harmonic means whose extended functions converge
in the full-field topology of the joint convergence and converge smoothly
on the buffer neighborhoods, in the interior finite-difference sense
used below.
\end{lemma}

\begin{proof}
The covariance limit is for the precise clamped form just defined.
The deterministic set \(C_{B,n}\) is optional for the cable field,
connected and nontrivial, and converges in Hausdorff distance to
\(\overline B\). The metric-local-set convergence theorem
\cite[Lemma~4.9]{ALS2018} therefore applies under the stated domain
hypotheses. It uses the ambient harmonic extensions of the revealed
and residual fields. Every limit is the deterministic local-set
decomposition at \(\overline B\), so the residual limit is the
zero-extended GFF on \(D\setminus\overline B\). Its finite pairing
vectors and their prelimit vectors are centered Gaussian. Convergence
of their characteristic functions forces convergence of their covariance
matrices, proving the displayed limit. This argument retains the whole
zero cable trace; it makes no assertion that deleting perimeter edges
commutes with harmonic extension.

Let \(\mathbf f=(f_i)_{i\le\ell}\) be a tuple of smooth tests, put
\(F_n=(\widehat\Phi_n(f_i))_{i\le\ell}\) and
\(F=(\Phi(f_i))_{i\le\ell}\), and let
\(Q_n=\E[F_n\mid\mathcal G_n(B)]\).
Linearity of the ambient extension and the cable Markov decomposition
give the exact identities
\begin{equation}
 \Cov(F_n-Q_n)=\mathsf K_n^B(\mathbf f,\mathbf f),\qquad
 \Cov(Q_n)=\widehat G_n^D(\mathbf f,\mathbf f)
                  -\mathsf K_n^B(\mathbf f,\mathbf f),
 \label{eq:projection-variance}
\end{equation}
where \(\widehat G_n^D\) is the covariance form of
\(\widehat\Phi_n\). The residual is independent of
\(\mathcal G_n(B)\), hence of \(T_n\). Both matrices in
\eqref{eq:projection-variance} converge. In the continuum, conditioning
on \(\mathcal G(B)\) has the corresponding killed covariance:
the functions in \(H^1_0(D)\) that vanish on \(B\) are the zero
extensions of \(H^1_0(D\setminus\overline B)\). Polygonal regularity
makes this the usual domain Markov decomposition.

Choose a countable dense family \((f_i^B)\) of interior smooth pairings
for the Gaussian Hilbert space \(\mathcal G(B)\), discarding dependent
coordinates in each finite list. Every cell meeting the support of a
fixed \(f_i^B\) belongs to \(B_n^\square\) for all sufficiently
large \(n\). Thus its ALS face pairing is exactly
\(\mathcal G_n(B)\)-measurable. Let \(Q_n^{(r)}\) be the Gaussian
projection of \(F_n\) onto the first \(r\) such pairings. For fixed \(r\),
convergence of the covariance matrices gives joint convergence with the
field to the corresponding finite projection \(Q^{(r)}\). If
\(Q=\E[F\mid\mathcal G(B)]\), density gives \(Q^{(r)}\to Q\) in
\(L^2\). Orthogonality and \eqref{eq:projection-variance} yield
\begin{equation}
 \lim_{r\to\infty}\limsup_{n\to\infty}
 \E|Q_n-Q_n^{(r)}|^2
 =\lim_{r\to\infty}\E|Q-Q^{(r)}|^2=0.
 \label{eq:projection-strong}
\end{equation}
This convergence of projection norms is the information needed for
an arbitrary local test.

For a bounded Lipschitz function \(v\) of \(F_n\), Gaussian conditioning
writes
\(\E[v(F_n)\mid\mathcal G_n(B)]=\Psi_n(Q_n)\), where \(\Psi_n\)
convolves \(v\) with the residual Gaussian law. Its covariance limit
and \eqref{eq:projection-strong} approximate this expression in \(L^1\),
uniformly for fine mesh, by bounded continuous functions of finitely
many interior pairings. The limit is
\(\E[v(F)\mid\mathcal G(B)]\). Multiply this identity by \(T_n\) and
by a bounded continuous function of finitely many interior pairings.
The \(L^1\) error remains valid because \(T_n\) is bounded. Passing to
the two limits proves the conditional factorization with \(T\) for every
full-field Gaussian cylinder given \(\mathcal G(B)\). A monotone-class
argument proves \eqref{eq:kernel-localization}. The vector and harmonic
mean versions follow by the same argument.
\end{proof}

\subsection{A domain-independent local kernel}

\begin{lemma}[Common local kernel]
\label{lem:local-kernel}
Fix \(K\Subset B\Subset P\) as above and a sequence of
\([0,1]\)-valued cable tests \(T_n\) supported in \(C_{K,n}\).
There are a subsequence and a measurable \(k\), with \(0\le k\le1\),
on continuum field restrictions to \(B\), such that:
\begin{enumerate}
\item In the zero-boundary reference field in \(P\), every joint limit
along this subsequence satisfies
\(\E[T\mid\Phi]=k(\Phi|_B)\).
\item The same \(k\) is the conditional kernel in every deterministic
domain \(D\) with \(\overline P\subset D\), for the same cable tests
in admissible approximations containing the identical subdivided
\(\mathsf P_n\). For a prescribed harmonic mean, including a constant
boundary value, require its ALS extended means to converge in
\(H^s(\mathcal B)\) to the chosen full continuum ALS extension, and to converge
smoothly on buffer neighborhoods in the interior finite-difference
sense used below.
Equivalently, for every bounded continuous field functional \(v\),
\[
 \E[T_n v(\widehat\Phi_n)]
       \longrightarrow \E[k(\Phi|_B)v(\Phi)]
\]
along the selected subsequence. This convergence of mixed expectations
does not require convergence of the tests as field functions.
\item Let a fixed random parameter \(\Theta\) select \(D\) and its mean,
and suppose that conditionally on \(\Theta\) the original total parent field
\(\Phi_D\) has that GFF law. Sample auxiliary cable fields conditionally
on \(\Theta\) in admissible approximations for almost every realization,
with the mean convergence required in item 2 for almost every realization.
For every joint limit
\[
 (\Theta,\widehat\Phi_n,T_n)
       \ \Longrightarrow\ (\Theta,\Phi_D^{\mathrm{aux}},T)
\]
along the selected subsequence, the auxiliary field marginal satisfies
\((\Theta,\Phi_D^{\mathrm{aux}})\stackrel{\mathrm d}=
  (\Theta,\Phi_D)\), and
\[
 \ind_{\{\overline P\subset D\}}
 \bigl(\E[T\mid\Theta,\Phi_D^{\mathrm{aux}}]
                  -k(\Phi_D^{\mathrm{aux}}|_B)\bigr)=0.
\]
\end{enumerate}
One subsequence works for any countable collection of tests, polygons,
and thresholds. For the collar tests at \(t\in a\Z\), write the selected
kernels as \(k_t^\pm\). In the restriction to \(B\) of every admissible
zero-boundary field \(g^D\),
\begin{equation}
       k_t^\pm(t+g^D|_B)=k_0^\pm(g^D|_B)
       \quad\hbox{almost surely}.
 \label{eq:kernel-threshold}
\end{equation}
The same identity holds under the conditional zero-boundary law of a
random selected domain. It is an identity under these Gaussian laws,
with the countable thresholds fixed in advance.
\end{lemma}

\begin{proof}
For the reference field, take a joint subsequential limit of
\((\widehat g_n^P,T_n)\). The reverse buffer lemma makes its conditional
mean a function \(k(g^P|_B)\). We show that this selected function also
governs the mixed expectations in other domains.

First translate the reference cable field by the linear cable
interpolation \(u_n\) of a fixed \(u\in C_c^\infty(P)\). Write
\(X_n=(\widetilde g_n^P,u_n)_{\nabla,n}\) for the Gaussian
Cameron--Martin coordinate; this is the isonormal pairing, not a
Dirichlet integral of a rough cable sample. The likelihood is
\begin{equation}
 L_{n,u_n}
 =\exp\left\{X_n-\tfrac12\|u_n\|_{\nabla,n}^2\right\}.
 \label{eq:kernel-cm}
\end{equation}
We need strong identification of \(X_n\), since \(T_n\) is arbitrary.
Let \(f=-\Delta u/(2\pi)\), and let \(\widehat u_n\) be the same
harmonic face extension of \(u_n\). The Gaussian reproducing identity
and linearity of the extension give
\[
 \Var X_n=\|u_n\|_{\nabla,n}^2,\qquad
 \Cov(X_n,\widehat g_n^P(f))=\int \widehat u_n f,\qquad
 \Var(\widehat g_n^P(f))=\widehat G_n^P(f,f).
\]
The three expressions tend respectively to
\(\|u\|_\nabla^2,\int uf=\|u\|_\nabla^2\), and
\(G_P(f,f)=\|u\|_\nabla^2\). Here the energy is a Riemann sum,
\(\widehat u_n\to u\) uniformly on compact subsets, and the last
limit is Green convergence. Hence
\begin{equation}
       \E|X_n-\widehat g_n^P(-\Delta u/(2\pi))|^2\longrightarrow0.
 \label{eq:kernel-cm-coordinate}
\end{equation}
The likelihood coordinate therefore converges jointly with every
bounded \(T_n\), not only marginally. For \(r>1\),
\begin{equation}
 \E[L_{n,u_n}^{\,r}]
   =\exp\left\{\tfrac12r(r-1)\|u_n\|_{\nabla,n}^2\right\}.
 \label{eq:kernel-cm-ui}
\end{equation}
These bounds give uniform integrability on bounded energy sets.
Tilting the reference joint law by \eqref{eq:kernel-cm} is exactly
the law of the translated total field and the test evaluated on that
translated cable field. The limiting likelihood is a function of the
continuum field. Such a tilt preserves its conditional kernel \(k\),
evaluated at the total restriction.

For two shifts in one reference law, the Gaussian relative entropy
bound and Pinsker's inequality give
\begin{equation}
       \|L_u-L_v\|_{L^1}
       \le \min\{2,\|u-v\|_\nabla\}.
 \label{eq:kernel-cm-continuity}
\end{equation}
The same inequality holds on the cables. Finite smooth nets in the
Dirichlet norm, together with \eqref{eq:kernel-cm-ui}, therefore make
the translated mixed-expectation limit uniform on compact families of
smooth shifts. This is equality of limiting measures tested first on
bounded continuous field functions. It does not require evaluating
a possibly discontinuous \(k\) on the discrete fields.

Now take a compared ambient graph containing the identical subdivided
\(\mathsf P_n\). Split at its boundary cut points and condition on
the outside cables and those boundary vertex values. The Markov
decomposition on \(\mathsf P_n\) is exactly
\[
              \widetilde\Phi_n=\widetilde g_n^P+H_n,
\]
where \(\widetilde g_n^P\) has the reference cable law
and is independent of \(H_n\). The mean \(H_n\) is discrete harmonic
at the interior vertices and linear on inside cables. Choose a smooth
cutoff \(\chi\) supported in \(P\) and equal to one near \(\overline B\).
At vertices set \(u_n=\chi H_n\) and interpolate linearly on cables.
It equals \(H_n\) on all cables needed for the pairings and tests in
\(B\). The reference and ambient harmonic extensions also agree on
those complete cells.

Interior discrete harmonic estimates and Green convergence make the
values and scaled finite differences of \(H_n\) converge in law on the
support of \(\chi\) to those of a smooth harmonic limit. They also
give convergence of the cutoff energies. This is a statement about
interior finite differences; the piecewise harmonic face extension
need not have a global classical smooth norm. Uniformly over the
domains, the covariance of \(H_n\) on a smaller buffer is bounded
in quadratic-form order by the harmonic covariance obtained from the
fixed outer box. Interior estimates consequently bound fixed smooth
norms of the limiting harmonic parts and their discrete analogues in
moments. Truncate such a bound, approximate the resulting family by a
finite smooth net in a weaker norm controlling the cutoff energy, and
apply \eqref{eq:kernel-cm-continuity}. Removing the truncation proves
the mixed-expectation limit on the restriction to \(B\). A prescribed
harmonic mean is treated by the same cutoff.

For completeness, the mean premise is automatic for the constant labels
used below in the regular inner class. For a constant \(t\), the full
continuum ALS extension \(u_t\) equals \(t\) on \(\overline D\cap\mathcal B\), is
harmonic on \(\mathcal B\setminus\overline D\), and has value zero on
\(\partial\mathcal B\). The prelimit mean equals \(t\) on
\(\Gamma_n\cap D_n\); its ambient extension is \(t\) times the planar
probability of hitting that cable domain before \(\partial\mathcal B\).
The maximum principle bounds its absolute value by \(|t|\).
For \(z\in\partial D\cap\mathcal B\), fix a radius \(r>0\) below its
distance to the box boundary and to a fixed interior point. For every
small \(\eta>0\), connectedness of \(D\) gives a compact path in \(D\)
from \(B(z,\eta)\) to outside \(B(z,r)\). Inner exhaustion supplies a
connected cable path with the same property, up to a vanishing mesh
error, for all sufficiently large \(n\). The planar Beurling estimate
\cite[Lemma~4.1]{ALS2018} bounds the probability of avoiding that path
until leaving \(B(z,r)\) by \(C(\eta/r)^\beta\). First let \(n\) tend to
infinity and then \(\eta\) to zero. The extended mean thus tends to \(t\)
at every such boundary point. Away from \(\partial D\), convergence is
\cite[Lemma~4.2(1)]{ALS2018}. Bounded convergence now gives convergence
to \(u_t\) in \(L^2(\mathcal B)\), hence in \(H^s(\mathcal B)\), even
when \(\partial D\) has positive area. On the buffer cells the mean is
exactly constant, so the interior finite-difference requirement also holds.

Thus the conditional kernel on that restriction is \(k\). Applying
\cref{lem:reverse-buffer} in the actual ambient domain promotes it to
conditioning on the full ambient field. The restricted Gaussian laws
are mutually absolutely continuous: conditional on \(H\), each is a
Cameron--Martin translate of the reference restriction with a strictly
positive likelihood, and integrating such translates preserves
equivalence. The representative of \(k\) is therefore unambiguous
under all these restrictions. The subsequence principle gives the
mixed-expectation convergence in item 2 even if a general bounded
test has several joint subsequential laws with the same conditional
mean.

For item 3, apply that deterministic mixed-expectation conclusion
for almost every fixed \(\Theta=\theta\). The full-field convergence
of the centered fields and the extended means identifies the limiting
conditional field marginal with the selected parent GFF law. Integrating
against bounded functions of \(\Theta\) gives
\((\Theta,\Phi_D^{\mathrm{aux}})\stackrel{\mathrm d}=
(\Theta,\Phi_D)\). The cutoff comparison identifies the conditional
mean given \((\Theta,\Phi_D^{\mathrm{aux}}|_B)\); then the reverse buffer
argument, applied conditionally after centering the selected mean,
identifies it given \((\Theta,\Phi_D^{\mathrm{aux}})\).
Integrate the bounded mixed expectations against bounded functions
of \(\Theta\), and use equality of measures and a monotone class.
This proves the displayed conditional identity on the auxiliary limit;
the displayed equality of marginals relates it to the separately given
original field. A countable choice of polygons is handled by partitioning
according to that choice.
The reference subsequence is never reselected for a realized domain.

Finally use a diagonal extraction for the countable tests and
thresholds. On \(C_{K,n}\) the exact identity
\(T_{n,t}^{\pm}(\widetilde g_n+t)=T_{n,0}^{\pm}(\widetilde g_n)\)
holds. The corresponding extensions differ by \(t\) on the complete
cells near \(B\), though they need not do so outside the selected
domain. Apply the two selected kernel statements to these identical
joint restricted laws. Their conditional means give
\eqref{eq:kernel-threshold}. The same mixed-expectation argument
conditional on \(\Theta\) gives its random-domain version.
\end{proof}

The cutoff estimate also gives the following fixed-buffer fact.

\begin{lemma}[Interior exhaustion]
\label{lem:kernel-exhaustion}
Let \(D_r\uparrow D\) be regular inner domains and \(B\Subset D_1\).
The laws of the zero-boundary continuum GFFs restricted to \(B\)
converge in total variation to the restriction of the GFF in \(D\).
This holds for simply connected \(D\) without a Jordan-boundary
assumption.
\end{lemma}

\begin{proof}
Choose a fixed polygon \(P\) with \(\overline B\subset P\Subset D_r\)
for all large \(r\). In the Markov coupling, the field in \(D\)
restricted to \(D_r\) is an independent zero-boundary field in \(D_r\)
plus a harmonic Gaussian \(H_r\), with covariance \(G_D-G_{D_r}\).
Interior Green convergence and harmonic estimates make \(H_r\to0\)
in mean square in each fixed smooth norm near \(\overline B\).
Cut it off inside \(P\). The cutoff Dirichlet norm tends to zero
in mean, and \eqref{eq:kernel-cm-continuity}, conditional on \(H_r\),
bounds the total variation distance by
\(\tfrac12\E[\min\{2,\|\chi H_r\|_\nabla\}]\). This tends to zero.
\end{proof}

\subsection{Forcing in the canonical parent}

We record precisely how the auxiliary test is used with the local-set
data of canonical ancestors, before revealing the parent remainder.
Let \(\Theta\) be that data, let \(\Phi_D\) be the
total field in the selected parent, and fix a collar with
\(\overline P\subset D\). Conditionally on \(\Theta\), sample the
auxiliary cable fields of \cref{lem:local-kernel}. For each deterministic
realization of the parent, the coupled convergence of excursion sets
\cite[Theorems~3 and~28]{ALS2023} and of FPSs
\cite[Proposition~4.7]{ALS2018} has the following consequence. In
every joint auxiliary limit along the already selected mesh subsequence,
the containing set is its canonical function of the auxiliary parent
field \(\Phi_D^{\mathrm{aux}}\). If
that set has the strict hits which force the chosen cable test, then
\(T=1\). For an excursion this assertion is made for each fixed
diameter rank; the countable union of these assertions handles an
excursion whose rank is random.

Here is the finite-mesh implication. A connected excursion approximant
with one hit in \(\operatorname{int}Q\) and another in
\(K\setminus\overline V\) has a sign path crossing the collar.
Along a path between the hits, an exterior point of \(V\) is reached
before any exit from \(K\); take that initial portion. For an FPS,
a point in \(\operatorname{int}Q\) is joined to the artificial parent
boundary by its prescribed one-sign path. For fine mesh the parent
contains the fixed \(\mathsf P_n\), so that path must exit \(V\) while
still inside \(K\), regardless of where the rest of the parent boundary
lies. Strict hits persist under compact-set convergence and give the
claimed implication. No pointwise value is assigned to the continuum
distribution.

For clarity, this conditional calibration does not require a measurable
family of Skorokhod representations over all possible parents. For a
fixed deterministic parent, take an arbitrary joint subsequence and
apply the cited coupled theorem and Skorokhod representation there.
A continuous hit-depth function
\(J_U(C)=\sup_{x\in C}\min\{1,\dist(x,U^c)\}\), with \(J_U(\varnothing)=0\), positive exactly
when \(C\) meets the strict open neighborhood \(U\), shows that the
product of the required hit depths with \(1-T\) is zero in every limit.
The resulting bounded
mixed-expectation conclusion holds along the selected subsequence.
Apply it pointwise to the measurable conditional auxiliary kernels
for \(\Theta=\theta\), and integrate by dominated convergence.
Together with item 3 of \cref{lem:local-kernel}, this gives
\[
       \E[T\mid\Theta,\Phi_D^{\mathrm{aux}}]
                      =k(\Phi_D^{\mathrm{aux}}|_B).
\]
If \(E^{\mathrm{aux}}\) is the corresponding canonical strict-hit event,
then \(E^{\mathrm{aux}}\in\sigma(\Theta,\Phi_D^{\mathrm{aux}})\) and
\(T=1\) on \(E^{\mathrm{aux}}\). Positivity gives
\(\ind_{E^{\mathrm{aux}}}(1-k(\Phi_D^{\mathrm{aux}}|_B))=0\).
The equality of the \((\Theta,\Phi_D^{\mathrm{aux}})\) and
\((\Theta,\Phi_D)\) marginals transfers this identity to the corresponding
canonical event \(E\) in the original parent:
\begin{equation}
             \ind_E\bigl(1-k(\Phi_D|_B)\bigr)=0
             \quad\hbox{almost surely}.
 \label{eq:kernel-forcing}
\end{equation}
Canonical ancestor data are
functions of the original field, so the identity remains true in
any extension carrying the stopped exploration. The later upper
bound will condition only on that stop; it does not assert a formula
for \(\E[T\mid\Phi_D,\text{stop}]\).

\begin{figure}[t]
\centering
\begin{tikzpicture}[font=\small,>=Latex]
 \draw[rounded corners,fill=gray!8] (0,0) rectangle (3.6,2.6);
 \draw[rounded corners,fill=teal!9] (.25,.25) rectangle (3.35,2.35);
 \draw[rounded corners,fill=white] (.55,.55) rectangle (3.05,2.05);
 \draw[dashed,rounded corners] (.85,.8) rectangle (2.25,1.8);
 \draw[fill=teal!20] (1.02,1.08) rectangle (1.38,1.48);
 \draw[thick,red,->] (1.2,1.28) .. controls (1.65,1.65) and (2.05,1.1) .. (2.7,1.35);
 \node at (1.2,.94) {\(Q\)};
 \node at (2.05,1.65) {\(V\)};
 \node[anchor=east] at (3.0,1.87) {\(K\)};
 \node[anchor=east] at (3.3,2.17) {\(B\)};
 \node[anchor=east] at (3.55,2.43) {\(P\)};
 \draw[->,thick] (3.85,1.3)--(4.55,1.3);
 \node[draw,rounded corners,align=center,text width=4.8cm] (ker) at (7.15,1.3)
   {one local function \(k(\Phi|_B)\)};
 \node[draw,rounded corners,align=center,text width=4.8cm] (force) at (2.5,-1.25)
   {Canonical parent and strict hits \(E\)\\\(\ind_E(1-k)=0\)};
 \node[draw,rounded corners,align=center,text width=4.8cm] (bound) at (8,-1.25)
   {Condition on the retained stop\\\(\Phi|_O=m+g^O\), collar means \(\to0\)};
 \draw[->] (ker.south west)--(force.north);
 \draw[->] (ker.south east)--(bound.north);
\end{tikzpicture}
\caption{The path test crosses \(V\) inside \(K\), and the separated
buffers \(B\) and \(P\) make its kernel common to different Gaussian
domains. Calibration in the canonical parent gives an identity on the
field space. The upper bound uses the conditional GFF law of the stopped
component.}
\label{fig:kernel-conditioning}
\end{figure}

\subsection{The zero-boundary collar bound}

Fix a bounded simply connected \(O\), a compact rational polygon
\(Q\Subset O\), and a regular inner exhaustion
\(Q\Subset O_1\Subset O_2\Subset\cdots\uparrow O\).
For the \(j\)-th collar take
\[
       V=O_j,\qquad K=O_{j+2},\qquad B=O_{j+3},\qquad P=O_{j+4}.
\]
Use regular inner cable approximations of \(O\) on the common grid.
They have complement-Hausdorff convergence, uniformly bounded
connectivity, and contain every fixed \(\mathsf P_n\) for fine mesh.
Write \(T_{j,n}^{\pm}\) for the zero-threshold tests. We claim
\begin{equation}
       \lim_{j\to\infty}\limsup_{n\to\infty}
                 \P(T_{j,n}^{\pm}=1)=0.
 \label{eq:zero-boundary-collar}
\end{equation}

A one-sign path in an interior buffer is contained in the closure of
one cable sign component. Interior vertex values are almost surely
nonzero. On an edge, the nonnegative or nonpositive intervals are
closures of Brownian excursions; allowing their zero endpoints does
not join distinct interior sign components. Thus a crossing component
has a fixed positive lower diameter once the collar is sufficiently
far from \(Q\).

We make the rank and limit order explicit. Let
\(\mathcal C_{r,n}\) be the cable excursion supports in decreasing
diameter order, with the ranking convention of \cite{ALS2023}, and
let \(\mathcal C_r\) be their continuum limits. Theorems 3 and 28
of that paper give joint Hausdorff convergence for every fixed finite
list of ranks. Theorem 1 gives finitely many positive-diameter supports
above any fixed size, and its CLE/FPS construction in Section 3 places
each support compactly inside \(O\). These facts hold for the rough
simply connected \(O\) and the regular complement-Hausdorff
approximations just chosen.

Fix \(j_0\) and slightly enlarge \(Q\) within \(O_{j_0}\). Let
\(d>0\) be the distance of this enlargement from
\(O\setminus O_{j_0}\). For \(j\ge j_0\) and fine mesh, every crossing
component has diameter at least \(d\). If one has rank larger than \(N\),
then \(\diam\mathcal C_{N+1,n}\ge d\). Fixed-rank convergence gives
\[
 \limsup_{n\to\infty}
 \P(\text{a crossing component has rank}>N)
 \le \P(\diam\mathcal C_{N+1}\ge d/2).
\]
The right side tends to zero as \(N\to\infty\) by positive-diameter
finiteness. For the first \(N\) limiting supports, compact containment
implies that all lie in \(O_{j-1}\) with probability tending to one as
\(j\to\infty\). At fixed \(j\), their Hausdorff convergence then puts
the first \(N\) approximants inside \(O_j\), since
\(\overline O_{j-1}\Subset O_j\). Consequently
\[
 \limsup_{j\to\infty}\limsup_{n\to\infty}
       \P(T_{j,n}^{\pm}=1)
 \le \P(\diam\mathcal C_{N+1}\ge d/2).
\]
Send \(N\to\infty\). This proves \eqref{eq:zero-boundary-collar}:
the mesh limit is first for each fixed collar, the collar moves next,
and the finite rank truncation is removed last.

Let \(k_{j,0}^{\pm}\) be the common kernels selected for these
countably many collar choices. For fixed \(j\), the mixed-expectation
conclusion of \cref{lem:local-kernel} in \(O\) identifies its integral
with the mesh limit of the bounded test. Thus
\begin{equation}
       \lim_{j\to\infty}
       \E[k_{j,0}^{\pm}(g^O|_{B_j})]=0,
       \qquad g^O\text{ a zero-boundary GFF in }O.
 \label{eq:zero-boundary-kernel}
\end{equation}
This proof applies directly to rough \(O\). The total variation statement
in \cref{lem:kernel-exhaustion} concerns a fixed buffer and supplies no
estimate with a collar moving toward the boundary.

\subsection{No crossing of a stopped component}

\begin{proposition}[No crossing]
\label{prop:no-crossing}
Let \(O\) be a component of a stopped exploration from
\cref{prop:exploration-thin}, with retained sigma-field \(\mathscr E\).
Suppose \(O\) lies in an initial free parent \(P_0\) of label \(b_0\in a\Z\), or in
one pointed wired parent \(P_0\) of label \(b_0=\pm a\). As provided there,
\(O\) is simply connected and, conditionally on \(\mathscr E\),
\[
              h|_O=m+g^O,\qquad m-b_0\in2a\Z.
\]
No canonical free cluster descended from that parent meets both \(O\)
and \(\partial O\). The conclusion holds simultaneously for all
pointed components and all recursive generations, and after a restart
in a captured odd layer with its new incoming label.

In the wired case \(O\) may be a residual component of the global
composite exploration of the original \(h\). The proof uses
\(O\subset P_0\), simple connectedness, and the
displayed conditional law. It requires no connectedness of the
explored set relative to that parent and no separately selected
lattice sweep inside a limiting wired hole.
\end{proposition}

\begin{proof}
First work in an indexed canonical parent \(D_{\mathrm{par}}\) whose
total field is \(b+f\).
For a positive first-generation cluster \(C\),
\cref{lem:canonical-geometry} supplies a containing positive excursion
at total level \(b\) and a containing upper FPS with sublevel paths at
total level \(b+2a\). For a negative cluster the containing paths are
below \(b\) and above \(b-2a\). The parent is the initial domain or a
canonical Jordan hole; its regular auxiliary approximations satisfy
the coupled convergence hypotheses used above.

Suppose \(C\) meets both \(O\) and \(\partial O\), and choose a rational
\(Q\Subset O\) whose interior it meets. In a regular inner exhaustion
as above, connectedness of \(C\) makes it meet \(\partial O_{j+1}\)
for every large \(j\). This gives the two strict neighborhoods for the
excursion test: the second lies in
\(O_{j+2}\setminus\overline O_j\). The FPS test needs only the strict
inner hit, because its boundary-connected path must exit the fixed
collar. On the relevant-parent event \(O\subset D_{\mathrm{par}}\),
all buffers are contained in this canonical parent.

The stopped label and this parent label belong to the same affine
lattice \(b_0+2a\Z\). For a positive cluster, either \(m\le b\), so
the excursion forces the above-\(m\) test, or \(m\ge b+2a\), so the
upper FPS forces the below-\(m\) test. For a negative cluster, either
\(m\ge b\), using the below-\(m\) excursion, or \(m\le b-2a\), using
the above-\(m\) FPS. These alternatives exhaust the lattice. The
absolute threshold \(m\) belongs to \(a\Z\).

Calibrate both tests at every fixed threshold in \(a\Z\) in the
canonical parent, conditional on its ancestor data. For the excursion,
use the forcing identity for each fixed containing rank and take their
countable union; for the FPS use the whole-set identity. The path
implications and \eqref{eq:kernel-forcing} then hold simultaneously
for the countable choices of parent index, rank, threshold, and collar.
They are identities on the original field space and survive adjoining
the stop. To state the bound globally, fix an indexed candidate \(C\)
and a rational \(Q\), and select the collar from the stopped state by
the rule described below. Let \(\mathsf H_j\in\mathscr E\) record that
the selected pointed component and its \(j\)th collar are defined with
the required clearances inside \(O\). Let \(\mathsf L_j\) record that
the indexed candidate and its parent \(D_{\mathrm{par}}\) exist and that
\(\overline P_j\subset D_{\mathrm{par}}\). This latter event may depend
on canonical data. Declare it false when \(\mathsf H_j\) fails, use the
empty set for an absent candidate or component, and put
\[
 \mathsf C_Q=\{C\cap\operatorname{int}Q\ne\varnothing,\,
                     C\cap\partial O\ne\varnothing\}.
\]
The right side below is defined to be zero outside \(\mathsf H_j\).
Substituting the stopped-measurable \(m\) in the simultaneous forcing
identities gives the global pointwise bound
\begin{equation}
 \ind_{\mathsf H_j\cap\mathsf L_j\cap\mathsf C_Q}
 \le \ind_{\mathsf H_j}
       \bigl(k_{j,m}^{+}(h|_{B_j})+k_{j,m}^{-}(h|_{B_j})\bigr).
 \label{eq:crossing-kernel-bound}
\end{equation}
The direction selected by the unknown parent label \(b\) has disappeared
from this inequality. In particular \(b\) need not be measurable in
\(\mathscr E\).

Condition now only on \(\mathscr E\). The event \(\mathsf H_j\), the
collar, and \(m\) are stopped-measurable. The component law gives
\begin{align*}
 \Pr(\mathsf H_j\cap\mathsf L_j\cap\mathsf C_Q\mid\mathscr E)
 &\le \ind_{\mathsf H_j}\E\!\left[
   k_{j,m}^{+}(m+g^O|_{B_j})+
   k_{j,m}^{-}(m+g^O|_{B_j})\,\middle|\,\mathscr E\right]
 \\&=
 \ind_{\mathsf H_j}\E\!\left[
   k_{j,0}^{+}(g^O|_{B_j})+
   k_{j,0}^{-}(g^O|_{B_j})\,\middle|\,\mathscr E\right],
\end{align*}
where \eqref{eq:kernel-threshold} is applied under that conditional
Gaussian law. For each realized \(O\), \eqref{eq:zero-boundary-kernel}
makes the right side tend to zero as \(j\to\infty\). On the branch where
this parent is relevant, a crossing with witness \(Q\) belongs to
\(\mathsf H_j\cap\mathsf L_j\) for every sufficiently large \(j\).
Conditional Fatou applied to the left indicators therefore gives zero
conditional probability for that relevant crossing. No conditioning
on \(\mathsf L_j\) is used.

All random choices can be made within the countable families already
selected. An open hit has a rational polygon \(Q\) as witness. The
closed-set coordinate of the stop determines whether a rational polygon
has positive clearance inside \(O\). Choose successively the first
regular rational polygon containing the preceding one and the next
finite list of rational compact buffers. Simple connectedness ensures
such a nested exhaustion exists, and the clearance tests make the choice
\(\mathscr E\)-measurable. Partition by each finite collar and threshold
choice before conditioning. The random-domain clause of
\cref{lem:local-kernel} and bounded convergence justify the displayed
conditional integrals. There are countably many components and
canonical clusters.

It remains to justify that the same argument reaches every recursive
parent. Assume it for the ancestors of a candidate cluster. If an
ancestor lies wholly in \(O\), its filled disk lies in \(O\): its
outer Jordan loop is inside the simply connected \(O\), so the
connected complement of \(O\) cannot enter that loop without crossing
it. All descendants then stay in \(O\). Otherwise, an ancestor that
does not cross \(O\) and does not lie inside it is disjoint from \(O\).
Connectedness puts \(O\) in one complementary canonical hole, unless
the candidate descendant is outside \(O\) already. In the relevant
hole the incoming label is some \(b\in b_0+2a\Z\). The same absolute
threshold argument applies there. On this branch \(O\subset D_{\mathrm{par}}\),
so \(\mathsf L_j\) holds for every sufficiently large collar and the
global bound applies. Calibration conditions on that
parent's canonical data, while the upper bound continues to condition
only on the original retained stop. Induction proves the claim.

For a global composite wired stop, \(A\) contains the wired layer, so
\(O\) lies in one pointed wired hole. Its label satisfies
\(m-b_0\in2a\Z\) with \(b_0=\pm a\), and all canonical descendants
have baselines in the same affine lattice. The preceding proof applies
in the original field using those absolute labels. No further
exploration is introduced inside the selected limiting hole.
\end{proof}
